% Add the listed directories to the search path
% (allows easy moving of files around later)
% these paths are searched AFTER local config kpsewhich

% *.sty, *.cls
\makeatletter
\def\input@path{{@resources/texlive/texmf-local/tex/latex/}
        ,{@resources/texlive/texmf-local/bibtex/bst/},
        ,{@resources/texlive/texmf-local/bibtex/bib/},
        ,{@local/}
        }
\makeatother
\makeatletter
\def\bibinput@path{{@resources/texlive/texmf-local/tex/latex/}
        ,{@resources/texlive/texmf-local/bibtex/bst/},
        ,{@resources/texlive/texmf-local/bibtex/bib/},
        ,{@local/}
        }
\makeatother
  % allow latex to find custom stuff
% this file is in the root directory so . is the path to the root
\providecommand{\econtexRoot}{}\renewcommand{\econtexRoot}{.}
                   % Set path to root dir
\documentclass{econark}
\usepackage{econark}                     % shared notation and shortcuts
\usepackage{llorracc-handouts}           % \handoutHeader, \handoutNameMake, \handoutC
\providecommand{\texname}{TractableBufferStock}% the handout this note belongs to: it shares its macros and references
\providecommand{\FigDir}{}\renewcommand{\FigDir}{\econtexRoot/Figures}
% The handout and its appendix link to each other as files in the same folder (this repo is the master copy):
% each PDF to the other PDF, and on the web, where make4ht defines \HCode, each page to the other page.
\ifdefined\HCode
  \providecommand{\HandoutPDF}{TractableBufferStock.html}
  \providecommand{\AppendixPDF}{TractableBufferStock-Appendix.html}
\else
  \providecommand{\HandoutPDF}{TractableBufferStock.pdf}
  \providecommand{\AppendixPDF}{TractableBufferStock-Appendix.pdf}
\fi
                 % macros for various subdirs
\usepackage{@local/\texname}             % macros specific to this handout
\usepackage{econark-bibfilesfind}        % Package to find valid bibfiles
\bibfilesfind{\texname}                  % find them
\usepackage{xr-hyper}                    % refer to the handout's labels, with links into its PDF
\usepackage{hyperref}                    % allow hyperlinks
\usepackage{econark-make4ht}             % link colors and pdf-or-web settings
\ifdefined\HCode% on the web, tex4ht takes no address: it links to the other page through that page's .xref
  \externaldocument[][nocite]{TractableBufferStock}% its labels, not its citations
\else
  \externaldocument[][nocite]{TractableBufferStock}[\HandoutPDF]% its labels, not its citations
\fi

% The book is the source of this text: its Supplemental Note content/supplemental/TractableBufferStockMacro.md,
% in intertemporal-choice.github.io. Make edits there first and carry them here; check-book-sync.py reports any
% differences.

\begin{document}
\handoutHeader

\begin{verbatimwrite}{\jobname.title}
A Macroeconomic Interpretation of the Tractable Buffer Stock Model
\end{verbatimwrite}

\handoutNameMake

Loosely following~\cite{cjSOE} (with some simplifications), this note extends the model of \href{\HandoutPDF}{A Tractable Model of Buffer Stock Saving} to analyze macroeconomic dynamics in a
small open economy with a large number of individuals, where the
population statistics reflect the fulfillment of individual consumers'
{\it ex ante} expectations; for example, exactly proportion $\urate$
of households who are employed in period $t$ become ``unemployed'' before
$t+1$, so that the aggregate labor supply of the ``active'' (still employed) members of a generation
evolves according to
\begin{equation}\begin{gathered}\begin{aligned}
  \PopnLvl_{t+1,t} & =  \erate \PopnLvl_{t,t}
,
\end{aligned}\end{gathered}\end{equation}
where the first subscript denotes the date being examined and the second denotes
the period of birth of the generation being examined.

We make strong assumptions that permit straightforward aggregation. The first such assumption is that newly unemployed households immediately migrate out of the
country (think of British retirees moving to southern
Spain).\footnote{The qualitative story is not changed if the unemployed stay at home and live off their savings; since they have a simple linear
  decision rule (they spend a constant proportion of their resources),
  accounting for their behavior is straightforward but
  complicates the exposition without adding much substance.  See \cite{cjSOE} for a model that incorporates a stay-at-home unemployed population.}
This means that macroeconomic variables will reflect only the circumstances of 
employed consumers, rather than a blend of the employed and the unemployed.

Each person is part of a single ``generation'' of households born at the
same time, and every new generation is larger by the factor $\PopnGroFac$
than the newborn generation in the previous period:
\begin{equation}\begin{gathered}\begin{aligned}
  \PopnLvl_{t+1,t+1} & =  \PopnGroFac \PopnLvl_{t,t}
.
\end{aligned}\end{gathered}\end{equation}

We assume that total production by the (surviving) members of a generation grows by
the factor $\WGroFac$ every period.  If total production is to grow
despite a shrinking number of surviving members of the generation, 
production {\it per active capita} must grow by $\WGroFac/\erate$, as the handout's \hyperref[eq:meanPreserve]{growth of individual labor productivity} provides.

Consider the economy in some period 0 in which the size of the newborn population
and the wage rate have been normalized to $\PopnLvl_{0,0} = \Wage_{0} = 1$, and index each generation by its age $\age$, the number of periods since its birth. If the economy was created $\bar{\age}$ periods before period 0, so that its oldest generation has age $\bar{\age}$, the ratio of the total population to the population
of newborns will be 
\begin{equation}\begin{gathered}\begin{aligned}
  1 + (\erate/\PopnGroFac)+ (\erate/\PopnGroFac)^{2}+ ... +(\erate/\PopnGroFac)^{\bar{\age}} & =  \left(\frac{1-(\erate/\PopnGroFac)^{\bar{\age}+1}}{1-(\erate/\PopnGroFac)}\right)
\end{aligned}\end{gathered}\end{equation}
whose limit as $\bar{\age} \rightarrow \infty$ is a finite number so long as $\erate/\PopnGroFac < 1$.  We assume that the population grows, $\PopnGroFac > 1$; since $\erate < 1$, that ensures both this and, below, a finite total labor income.

Relative to the labor income of period 0's newborn cohort ($\PopnLvl_{0,0} \Wage_{0} = 1$),
the total labor income in period 0 of the generation born in period $-1$ is $\PopnGroFac^{-1}$; the sum of the incomes of all of the two-period-old individuals is $\PopnGroFac^{-2}$, and so on; total labor income for all generations in the economy in period 0 is
\begin{equation}\begin{gathered}\begin{aligned}
  1 + \PopnGroFac^{-1}+ \PopnGroFac^{-2}+ ...+\PopnGroFac^{-\bar{\age}} & =  \left(\frac{1-(\PopnGroFac^{-1})^{\bar{\age}+1}}{1-\PopnGroFac^{-1}}\right) \label{eq:LabIncTot}
,
\end{aligned}\end{gathered}\end{equation}
which is finite because $\PopnGroFac>1$.  
The {\it proportion} of aggregate income accounted for by a generation
born at any specific moment declines toward zero as time passes (old generations never die,
they just fade away).

In the balanced growth equilibrium, the growth factor for aggregate population
will therefore be $\PopnGroFac$ and output per capita will increase by $\WGroFac$ per period.  Total
labor income therefore grows by $\PopnGroFac \WGroFac.$


\section{Stakes}\label{sec:TractableBufferStockMacro-Stakes}

We now examine this model under two assumptions about the initial
``stake'' of newborns in the economy.  (We use ``stake'' to designate a
transfer received by newborns).  This is explicitly {\it not} an
inheritance, as we have assumed that individuals have no bequest
motive and newborns are unrelated to anyone in the existing
population.  Our motivation is to make the model more tractable,
rather than to represent an important feature of the real world; the model without a ``stake'' can only be solved by simulation, and Figure~\ref{fig:SOENoStakescPath} shows how its aggregate consumption evolves when households start with no wealth.  Both figures are drawn by the handout's \href{https://github.com/llorracc/TractableBufferStock/blob/main/TractableBufferStock.ipynb}{notebook}; run it in your browser, with nothing to install, in \href{https://intertemporal-choice.github.io/lite/notebooks/index.html?path=TractableBufferStock.ipynb}{JupyterLite}, or open it in \href{https://colab.research.google.com/github/llorracc/TractableBufferStock/blob/main/TractableBufferStock.ipynb}{Colab} or \href{https://mybinder.org/v2/gh/llorracc/TractableBufferStock/main?labpath=TractableBufferStock.ipynb}{Binder}.

\subsection{A \texorpdfstring{``Stake''}{\textquotedblleft Stake\textquotedblright} That Yields a Representative Agent}\label{sec:TractableBufferStockMacro-StakeRepAgent}

We first consider a version of the model in which an exogenous
redistribution program guarantees that the behavior of employed households can 
be understood by analyzing the actions of a ``representative employed agent.''  

If a benevolent source outside the economy were to provide every newborn with an initial
transfer upon birth of size $\bTrgNrm^{\hideSup}$, then the newborn's total monetary resources 
would be
\begin{equation}\begin{gathered}\begin{aligned}
  \mNrmE_{t,t} & =  \bTrgNrm^{\hideSup} + 1
\\ & =  \mTrgNrm^{\hideSup} 
.
\end{aligned}\end{gathered}\end{equation}

Thus, per-capita market resources for members of the newborn generation
would be exactly equal to the target level of market resources for a
person anticipating the future path of labor income that the members of the newborn
generation actually anticipate (which is the same as the future path
anticipated by all other generations as well).

If such a transfer policy had been in place forever, the economy at every 
point in time would consist of employed households whose consumption had 
been equal to its steady-state value ${\cNrm}^{\EmpLbl}$ for their whole lives.  That is, every 
individual agent in this economy would be identical in their {\it ratio} of consumption,
market resources, etc.\ to permanent labor income.  
The behavior of any individual would therefore be fully captured by the behavior
of a representative employed agent.\footnote{The {\it level} of permanent 
labor income will differ for different households, depending on their age
and accumulated wage seniority; the circumstances of individuals are identical
only after their problem has been normalized by their varying levels of permanent 
labor income.}

The foregoing scenario assumed that the ``stake'' is provided by a mysterious
``benevolent source outside the economy.''  Fortunately, there is an
easy way to eliminate this problematic assumption: Assume that the
stakes are financed by a wage tax.

The size of the required tax rate is calculated as follows.  The total size
of the resources transferred to the newborn generation must be
\begin{equation}\begin{gathered}\begin{aligned}
  \check{\BLvl}^{\EmpLbl}_{t,t} & =  \check{\bNrm}^{\hideSup} \PopnLvl_{t,t} \check{\Wage}_{t}
\end{aligned}\end{gathered}\end{equation}
where
\begin{equation}\begin{gathered}\begin{aligned}
  \check{\Wage}_{t} & =  \left(1-\TaxRte\right)\Wage_{t}
\end{aligned}\end{gathered}\end{equation}
is the after-tax wage rate for the economy as a whole (and $\check{\bNrm}^{\hideSup}$ is the steady state target ratio of bank balances to after-tax wages).

From \eqref{eq:LabIncTot}, the ratio of total aggregate labor income 
to the labor income of the newborn generation is 
\begin{equation}\begin{gathered}\begin{aligned}
  \left(\frac{1}{1-\PopnGroFac^{-1}}\right)
\end{aligned}\end{gathered}\end{equation}
so, since the stake is $\check{\bNrm}^{\hideSup}$ times the after-tax wage, the aggregate wage tax rate required to finance a ``stake'' of size $\check{\bNrm}^{\hideSup}$ for
newborns is given by 
\begin{equation}\begin{gathered}\begin{aligned}
  \check{\bNrm}^{\hideSup}(1-\TaxRte) & =    \left(\frac{\TaxRte}{1-\PopnGroFac^{-1}}\right)
\\ \TaxRte & =  \left(\frac{(1-\PopnGroFac^{-1})\check{\bNrm}^{\hideSup}}{1+(1-\PopnGroFac^{-1})\check{\bNrm}^{\hideSup}}\right)
.
\end{aligned}\end{gathered}\end{equation}

Note, however, that in an economy where this tax has existed forever, the consequence
of the tax is effectively just a permanent reduction in after-tax labor income by proportion $\TaxRte$, compared to its value in the absence of the tax.  
Given the homotheticity of the model, a permanent rescaling by a constant 
factor leaves the scaled version of the individual's problem (and its solution)
unchanged.  Thus we can conclude not only that a representative agent exists
in this economy, but that the steady-state characteristics of the 
representative agent's problem are identical (in ratio form) to the 
characteristics of the unrescaled individual's problem; that is, 
$\check{\cFunc}(\mNrm) = \cFunc^{\EmpLbl}(\mNrm)$, $\check{\bNrm}^{\hideSup} = \bTrgNrm^{\hideSup}$, and so on.  

Matters are not much more complicated outside the balanced growth
steady state, so long as we assume that the government always
transfers the amount $\check{\bNrm}^{\hideSup}$ to newborn households, financed by
the tax $\TaxRte$ derived above.  Consider, for example, an economy that
was in steady-state equilibrium leading up to period $\prdt$, and at
the beginning of $\prdt$ there is a sudden realization that future
growth rates will be higher than those anticipated and experienced in
the past: $\Alt{\WGroFac} > \WGroFac$ after $\prdt$.  Since expected
growth rates affect $\bTrgNrm^{\hideSup}$, the tax rate must be immediately and
permanently changed so that the generations born after $\prdt-1$
receive a ``stake'' of the proper new size.  This change in $\TaxRte$ has
two consequences for the generations that survive from periods prior
to $\prdt$.  Under the old tax rate, they would have experienced
$\bNrmE_{t} = \bLvl_{t}/((1-\TaxRte)\Wage_{t}) = \bTrgNrm^{\hideSup}$; the change in
expectations has no effect on $\bLvl_{t}$ or $\Wage_{t}$ but changes
the tax rate to $\Alt{\TaxRte}$.  Thus these households will have an
actual resource ratio that differs from its new target value, $\bNrmE_{t}
\neq \Alt{\bTrgNrm}^{\hideSup}$, both because the after-tax income scaling factor
has changed and because the target ratio has changed from $\bTrgNrm^{\hideSup}$
to $\Alt{\bTrgNrm}^{\hideSup}$.

However, if we started out in steady-state, the {\it ratio} problem of every
member of the continuing-employed population is identical to that of every 
other such household (though, again, their masses differ depending on age, etc); as a result, the dynamics of the economy are fully captured
by keeping track of the relative weights in the economy of the (gradually diminishing) ``representative shocked agent'' and the (gradually increasing) ``representative new agent'' 
whose behavior is locked at its steady-state value.\footnote{
If the economy has experienced multiple shocks, it will be necessary to
retain a complete history of the shocks in order to compute the properly 
population-weighted dynamics.  This is not too hard to do, if we start with
the assumption that the economy started at its balanced growth equilibrium before the shocks began to arrive.}

Figure~\ref{fig:TractableBufferStockMacro-SOEStakescPath} illustrates the
dynamics in this economy using an experiment identical to the handout's \hyperref[sec:FinalExperiment]{final experiment} for the individual's problem: In period 0 there is a one-off
decline in the future discounting rate (assuming the economy was in
steady state before period 0).  In that experiment, each individual
consumer's consumption function shifted down, and consumption
experienced a discrete jump downward, because the agent became more
patient.  Here, there is a modest further effect: With
more-patient consumers, the tax rate that the government sets to finance a transfer
of $\bTrgNrm^{\hideSup}$ to the newborns must be larger (so that the ratio of initial assets to
after-tax income is larger).  The higher tax also lowers the after-tax income of the households alive at the change, so they start from $1+\bTrgNrm^{\hideSup}(1-\TaxRte)/(1-\Alt{\TaxRte})$, a little above the old target, as the figure shows.  Qualitatively, the
dynamics are indistinguishable from the individual consumer's dynamics
obtainable without working through the extra complication involved in
accounting for the ``stakes.''

\begin{figure}
\caption{Aggregate $\CNrm$ in a Small Open Economy Before and After $\discRte$ Decline}
\includegraphics[width=6in]{\FigDir/SOEStakescPathAfterThetaDropPlot}
\label{fig:TractableBufferStockMacro-SOEStakescPath}
\FigNote{With stakes, aggregate consumption behaves like an individual's: it drops, then rises to a higher ratio.}
\end{figure}

\subsection{No Stake}\label{sec:TractableBufferStockMacro-NoStake}

The polar alternative to assuming that newborns get a ``stake'' is to
assume that newborns enter the economy with zero assets.  Analysis of
this version of the model must be performed using simulation methods,
because households of different ages will have different levels of
assets.  (Because the consumption function is concave and has no closed form, aggregation cannot be done analytically.)

Our simulation procedure assumes that at date 0 the economy has existed forever (so 
that the age distribution of relative populations and productivities are at their steady-state values), but saving has been impossible prior to
period 0.\footnote{In periods before 0, unemployment presumably would have 
meant immediate death by starvation; we think of this more as a starting point
for the simulations than as a realistic description of a plausible economy.}
With everyone's $\bNrmE_{t}=0$, the ratio of market resources to permanent labor income is the same for all individuals, whatever their age $\age$:
\begin{equation}\begin{gathered}\begin{aligned}
  \mNrmE_{0,-\age} & =  1
.
\end{aligned}\end{gathered}\end{equation}
The consumption ratio in period 0 is therefore $\cFunc(1)$ for every
household (regardless of age), while the
level of per-capita labor income for a generation of age $\age$ is $\erate^{-\age}$.\footnote{The absolute level of wages will have
grown by $\WGroFac/\erate$ per period for these households since their birth, but we have normalized by the level of wages
for period-0 newborns, which cancels the $\WGroFac$ from the expression.}  The population of
such workers is $(\erate/\PopnGroFac)^{\age}$, so aggregate consumption will
be given by the per-capita consumption ratio, multiplied by the
per-capita level of permanent income, multiplied by the population of
workers still alive:
\begin{equation}\begin{gathered}\begin{aligned}
  \CLvl_{0} & =  \sum_{\age=0}^{\infty} \cFunc(1) \erate^{-\age} (\erate/\PopnGroFac)^{\age}
\\ & =  \cFunc(1) \sum_{\age=0}^{\infty}  \PopnGroFac^{-\age}
\\ & =  \cFunc(1) \left(\frac{1}{1-\PopnGroFac^{-1}}\right)
.
\end{aligned}\end{gathered}\end{equation}

The longer a generation lives, the more time it will have had to save toward
its target level of wealth; but newborns always begin life with no assets.  After period 0, therefore, age-heterogeneity in assets and consumption ratios creeps into the population.

The foregoing discussion contains (in some cases implicitly) all the
assumptions necessary to conduct a simulation of this economy.
Figure~\ref{fig:SOENoStakescPath} shows the path of $\CNrm_{t}$, aggregate consumption as a ratio to aggregate labor income, starting from period 0 for an economy under the benchmark parameterization of \href{\HandoutPDF}{A Tractable Model of Buffer Stock Saving}.  The only extra parameter required beyond the handout's is $\PopnGroFac$; we choose $\PopnGroFac=1.01$ corresponding
roughly to the postwar population growth rate in the United States.

\begin{figure}
\caption{Path of Aggregate $\CNrm$ in a Small Open Economy Without Stakes From Date 0}
\includegraphics[width=6in]{\FigDir/SOENoStakescPath}
\label{fig:SOENoStakescPath}
\FigNote{Starting with no wealth, aggregate consumption rises quickly toward its steady-state ratio as households build their buffer stocks.}
\end{figure}

\pagebreak
\bibliographystyle{econark}
\bibliography{\bibfilesfound}

\end{document}
